\subsection{SP4: juego distribuido para navegación y acoplamiento seguro}
\label{subsec:sp4-v4-metodo}

SP4 conecta la selección mecánica de SP3 con el contacto físico: una coalición ya
reclutada debe alcanzar sus puertos sin colisiones, sin exceder torque de rueda y sin
quedar bloqueada por decisiones simétricas. La contribución no es un planificador
multiagente monolítico. Es una composición de tres capas: (i) un juego poblacional
local de vivacidad, (ii) una clausura distribuida de derecho de paso y admisión por
colores, y (iii) un filtro HOCBF en el espacio de aceleraciones ejecutables. La
arquitectura adapta la lógica de juegos poblacionales distribuidos de
\citet{quijano2017population,martinezpiazuelo2022automatica} a robots móviles con
dinámica, obstáculos y pares de rueda.

\begin{figure}[H]
\centering
\begin{tikzpicture}[
  node distance=6mm and 7mm,
  block/.style={draw,rounded corners,align=center,minimum height=10mm,minimum width=27mm,fill=blue!6},
  gate/.style={draw,rounded corners,align=center,minimum height=10mm,minimum width=27mm,fill=orange!10},
  arr/.style={-{Latex[length=2mm]},thick}
]
\node[block] (astar) {A* y cola radial\\por puerto};
\node[block,right=of astar] (game) {Juego local\\$\{\mathrm{go,slow,wait}\}$};
\node[gate,right=of game] (cap) {Proyección de capacidad\\y dueño del recurso};
\node[gate,below=of cap] (hocbf) {HOCBF de pares,\\obstáculos y carga};
\node[block,left=of hocbf] (torque) {Proyección de\\torque y aceleración};
\node[block,left=of torque] (plant) {Uniciclo dinámico\\sin reparar posición};
\draw[arr] (astar)--(game);
\draw[arr] (game)--(cap);
\draw[arr] (cap)--(hocbf);
\draw[arr] (hocbf)--(torque);
\draw[arr] (torque)--(plant);
\draw[arr,dashed] (plant.south) -- ++(0,-7mm) -|
  node[pos=0.28,below,align=center,font=\scriptsize]{estado, conflictos y robots acoplados}
  (game.south);
\end{tikzpicture}
\caption{Cadena ejecutable de SP4 v4. La capa de juego decide vivacidad; la HOCBF
modifica aceleraciones, no posiciones.}
\label{fig:sp4-v4-chain}
\end{figure}

\subsubsection{Juego de recursos de pareja}

Cada robot $i$ mantiene una distribución
$x_i\in\Delta_3$ sobre $\mathcal A=\{g,s,w\}$, correspondientes a avanzar,
avanzar lento y esperar. Sea $\rho=(i,j)$ un recurso asociado a una arista activa del
grafo de conflicto $\mathcal G_c(t)$. Con intensidades
$q=(1,0{,}42,0)^\top$, su ocupación es
\begin{equation}
 z_\rho(x)=c_{i\rho}q^\top x_i+c_{j\rho}q^\top x_j,
 \qquad z_\rho(x)\leq 1,
 \label{eq:sp4-pair-capacity}
\end{equation}
donde $c_{i\rho}=1$ si el robot participa en el recurso. La función potencial usada en
cada grafo congelado es
\begin{equation}
 \Phi(x)=\sum_i c_i^\top x_i+
 \frac{\beta}{2}\lVert z(x)\rVert_2^2+
 \frac{\varepsilon}{2}\sum_i\lVert x_i\rVert_2^2,
 \qquad \beta=0{,}52,\quad\varepsilon=0{,}08.
 \label{eq:sp4-potential}
\end{equation}
El término lineal premia progreso y prioridad wrench; la edad de espera solo crece si el
robot está admitido y cede realmente, y se satura para que no domine el precio dual.
El gradiente local es
\begin{equation}
 \nabla_{x_i}\Phi=c_i+\beta\sum_\rho c_{i\rho}qz_\rho+
 \varepsilon x_i.
 \label{eq:sp4-gradient}
\end{equation}

Para protocolo primal--dual se actualizan preferencias y precios persistentes mediante
\begin{align}
 x_i^{\ell+1}&=\Pi_{\Delta_3}\!\left[x_i^\ell-\eta
 \left(\nabla_{x_i}\Phi+\sum_\rho c_{i\rho}q\lambda_\rho\right)\right],\\
 \lambda_\rho^{\ell+1}&=\left[\lambda_\rho^\ell+gamma
 (z_\rho-1)\right]_+.
 \label{eq:sp4-pd}
\end{align}
La variante replicator usa el paso espejo entrópico
\begin{equation}
 x_{ia}^{\ell+1}=\frac{x_{ia}^{\ell}
 \exp[-\eta g_{ia}^{\ell}]}{\sum_bx_{ib}^{\ell}
 \exp[-\eta g_{ib}^{\ell}]},
 \label{eq:sp4-replicator}
\end{equation}
con $g_i$ igual al gradiente aumentado. Tras cada revisión se aplica una proyección
local del recurso: si $z_{ij}>1$, ambos robots escalan proporcionalmente su masa
$\{g,s\}$ y trasladan el exceso a $w$. Como esta operación solo reduce actividad, una
restricción ya satisfecha no vuelve a violarse durante la barrida.

Para un grafo estático, el conjunto factible es compacto y convexo. Como
$\varepsilon>0$, \eqref{eq:sp4-potential} es estrictamente convexa; por tanto, posee un
minimizador único y sus condiciones KKT caracterizan el equilibrio variacional del
juego restringido. La convergencia asintótica de \eqref{eq:sp4-pd} requiere los
supuestos habituales sobre pasos y grafo estático. En la implementación híbrida, con
grafo, rutas y admisión variables, no se reclama esa convergencia global: se auditan
residual KKT, simplex y capacidad en cada ejecución. Esta separación evita convertir
un resultado condicional en un teorema sobre todo el sistema.

La preferencia continua se convierte en escala de velocidad mediante
$\sigma_i=q^\top x_i$. Si dos predicciones comparten recurso, un protocolo local conserva
un dueño hasta que desaparece el conflicto; el dueño recibe
$\sigma_i\geq0{,}8$ y el otro cede. En escenarios densos, una coloración conservadora
admite una clase por fase. Para el paso estrecho se priorizan primero los puertos con
menor clearance geométrico, evitando que contactos anteriores cierren el acceso.

\subsubsection{HOCBF y límites de rueda}

El estado dinámico del robot es
$s_i=(p_i,\theta_i,v_i,\omega_i)$, con
$\dot p_i=v_ie(\theta_i)$, $\dot v_i=a_i$ y
$\dot\omega_i=\alpha_i$. Para dos discos se define
\begin{equation}
 h_{ij}=\lVert p_i-p_j\rVert_2^2-d_{ij}^2.
 \label{eq:sp4-hij}
\end{equation}
La restricción HOCBF de orden relativo dos es
\begin{equation}
 \ddot h_{ij}+(k_1+k_2)\dot h_{ij}+k_1k_2h_{ij}\geq0,
 \qquad k_1=k_2=2{,}8,
 \label{eq:sp4-hocbf}
\end{equation}
y es afín en las aceleraciones longitudinales $a_i,a_j$. Se construye análogamente
para obstáculos circulares y la carga. El giro nominal conduce la ruta; $\alpha_i$ no
aparece directamente en \eqref{eq:sp4-hocbf} para el centro del uniciclo, una limitación
que motiva márgenes A* y replanificación en cuellos. La formulación sigue el principio
de filtrado de seguridad de \citet{ames2017cbf}, pero se aplica aquí después del juego y
antes de los actuadores.

La conversión a ruedas es
\begin{equation}
 \tau_{R,L}=r_w\left(\frac{m_ia_i}{2}
 \pm\frac{I_i\alpha_i}{L_i}\right),
 \qquad |\tau_{R,L}|\leq\bar\tau_i.
 \label{eq:sp4-wheel-torque}
\end{equation}
Las cajas de $a_i,\alpha_i,v_i,\omega_i$ y los cuatro semiespacios de torque se incluyen
en la proyección alternada. Después se saturan físicamente los torques, se reconstruye
la aceleración realizada y se vuelve a evaluar \eqref{eq:sp4-hocbf}. No se corrigen
posiciones después de integrar.

La invariancia continua de \eqref{eq:sp4-hocbf} es condicional a factibilidad e
inicialización compatible. No implica automáticamente invariancia de la integración
discreta. Por ello se reportan por separado residual ejecutado y clearance barrido; el
confirmatorio contiene un contraejemplo de baseline donde residual casi nulo coexiste
con una penetración discreta de $5{,}7$ mm.

\subsubsection{Complejidad y protocolo experimental}

Con $A=3$, $I_g$ revisiones y $|\mathcal E_c|$ conflictos, el juego local cuesta
$\mathcal O(I_gA|\mathcal E_c|)$ y comunica dos mensajes por arista y revisión. La
implementación secuencial HOCBF cuesta
$\mathcal O(I_b(N^2+NM))$ para $M$ obstáculos. La referencia central enumera subconjuntos
y, en esta implementación, crece como $2^N$; no se usa más allá de flotas pequeñas.
Así, el juego continuo de SP4 no es el componente NP-hard: la dificultad combinatoria
queda en la referencia central y en el planeamiento multi-robot conjunto, que aquí se
descompone en A*, coloración y filtros locales.

Se comparan cinco métodos bajo mundos pareados: directo+HOCBF, prioridad local+HOCBF,
referencia central+HOCBF, replicator distribuido+HOCBF y primal--dual distribuido+HOCBF.
No se etiqueta ningún proxy como ORCA: ORCA real no está implementado en esta campaña.
El desarrollo usa una semilla; el confirmatorio congela parámetros y usa dos semillas
nuevas, seis escenarios y $N=4$ (60 ejecuciones). Un certificado de stress separado
prueba el método propuesto con $N=12$; no sustituye potencia estadística multi-semilla.

Como comprobación visual posterior se construye una única escena pareada en el motor
real de CoppeliaSim: dos copias de \texttt{narrow\_passage}, semilla 886001, muestran
respectivamente directo+HOCBF y primal--dual+HOCBF. Las trayectorias ejecutables se
embeben en el archivo \texttt{.ttt} mediante Lua y se reproducen con paso $0{,}12$ s.
El gate exige escena guardada, capturas reales, error de reproducción inferior a 1 mm y
reproducción de ambos desenlaces. Esta prueba valida geometría, reproducibilidad y
plausibilidad visual; no es un lazo dinámico independiente ni evidencia de hardware.